\chapter{Quadratic Forms and Definiteness}

A quadratic form is a homogeneous polynomial of degree $2$ in several
variables. In economics it shows up as the variance of a portfolio,
$w^{T} \Sigma w$, and as the second-order term in the Taylor expansion of an
objective function. Whether the quadratic form is positive, negative, or
mixed in sign decides whether a critical point of a smooth function is a
(local) minimum, maximum, or saddle.

\section{Core Rules}

\begin{rulebox}{Quadratic form}
Given a symmetric matrix $A \in \R^{n \times n}$, the associated quadratic
form is the function
\[
Q(x) \;=\; x^{T} A x \;=\; \sum_{i, j} a_{ij} x_{i} x_{j}, \qquad x \in \R^{n}.
\]
Without loss of generality one may always take $A$ symmetric: replace any
$A$ by $\tfrac{1}{2}(A + A^{T})$ without changing $Q$.
\end{rulebox}

\begin{rulebox}{Definiteness}
Let $A \in \R^{n \times n}$ be symmetric. We say:
\begin{itemize}[leftmargin=1.6em,itemsep=0.3ex]
    \item $A$ is \emph{positive definite} (pd) if $x^{T} A x > 0$ for all $x \neq 0$.
    \item $A$ is \emph{positive semidefinite} (psd) if $x^{T} A x \ge 0$ for all $x$.
    \item $A$ is \emph{negative definite} (nd) if $x^{T} A x < 0$ for all $x \neq 0$.
    \item $A$ is \emph{negative semidefinite} (nsd) if $x^{T} A x \le 0$ for all $x$.
    \item $A$ is \emph{indefinite} if $x^{T} A x$ takes both positive and negative values.
\end{itemize}
\end{rulebox}

\begin{rulebox}{Eigenvalue criterion}
For symmetric $A$:
\begin{itemize}[leftmargin=1.6em,itemsep=0.3ex]
    \item $A$ pd $\Leftrightarrow$ all eigenvalues $> 0$;
    \item $A$ nd $\Leftrightarrow$ all eigenvalues $< 0$;
    \item $A$ psd $\Leftrightarrow$ all eigenvalues $\ge 0$;
    \item $A$ nsd $\Leftrightarrow$ all eigenvalues $\le 0$;
    \item $A$ indefinite $\Leftrightarrow$ $A$ has both a positive and a negative eigenvalue.
\end{itemize}
\end{rulebox}

\begin{rulebox}{Principal-minor (Sylvester) criterion}
Let $D_{k}$ denote the determinant of the leading $k \times k$ principal
submatrix of $A$ (the top-left $k \times k$ block). Then:
\begin{itemize}[leftmargin=1.6em,itemsep=0.3ex]
    \item $A$ pd $\Leftrightarrow$ $D_{1}, D_{2}, \dots, D_{n} > 0$;
    \item $A$ nd $\Leftrightarrow$ $(-1)^{k} D_{k} > 0$ for $k = 1, \dots, n$
          (signs alternate: $D_{1} < 0, D_{2} > 0, D_{3} < 0, \dots$).
\end{itemize}
For semidefiniteness one must check \emph{all} principal minors, not just
the leading ones; the simple form above does not extend.
\end{rulebox}

\begin{rulebox}{Leading minors: the other sign patterns}
Suppose \emph{every} leading minor is non-zero, $D_{1}, D_{2}, \dots, D_{n}
\neq 0$. Then $A$ is invertible and exactly one of three things holds:
\begin{itemize}[leftmargin=1.6em,itemsep=0.3ex]
    \item all $D_{k} > 0$: positive definite;
    \item signs $-, +, -, \dots$ (that is, $(-1)^{k} D_{k} > 0$): negative definite;
    \item \textbf{any other pattern}: indefinite. For $n = 3$, for example,
          $D_{1} < 0,\ D_{2} > 0,\ D_{3} > 0$ or $D_{1} > 0,\ D_{2} > 0,\ D_{3} < 0$.
\end{itemize}
A negative principal minor of \emph{even} order settles the question at
once, even when other minors vanish: psd and nsd both need every
even-order principal minor to be $\ge 0$. For example, if $D_{2} < 0$ (or
any $2 \times 2$ principal minor is negative), the form is already
indefinite on the plane of those two coordinates, so $A$ is indefinite
whatever the remaining minors are.
\end{rulebox}

\begin{rulebox}{Diagonal shortcuts}
The diagonal entry $a_{ii} = e_{i}^{T} A e_{i}$ is the value of the form at
the $i$-th unit vector. Hence, before computing any determinant:
\begin{itemize}[leftmargin=1.6em,itemsep=0.3ex]
    \item pd needs \emph{all} $a_{ii} > 0$, nd needs all $a_{ii} < 0$
          (psd needs all $a_{ii} \ge 0$, nsd all $a_{ii} \le 0$);
    \item two diagonal entries of \textbf{opposite signs} $\Rightarrow$ indefinite;
    \item a \textbf{zero} diagonal entry $a_{ii} = 0$ with some $a_{ij} \neq 0$
          in its row $\Rightarrow$ indefinite, because the $2 \times 2$
          principal minor on rows and columns $i, j$ equals
          $0 \cdot a_{jj} - a_{ij}^{2} < 0$;
    \item a diagonal matrix: the diagonal entries are the eigenvalues, read
          the verdict off their signs.
\end{itemize}
The conditions in the first bullet are only necessary: all $a_{ii} > 0$
does \emph{not} make $A$ positive definite.
\end{rulebox}

\begin{rulebox}{Semidefiniteness: all seven principal minors ($3 \times 3$)}
When some leading minor is $0$, Sylvester's test stops. For a symmetric
$3 \times 3$ matrix list \emph{all} principal minors (delete the same rows
and columns):
\[
\underbrace{a_{11},\ a_{22},\ a_{33}}_{\text{order }1},\qquad
\underbrace{M_{12},\ M_{13},\ M_{23}}_{\text{order }2},\qquad
\underbrace{\det A}_{\text{order }3},
\]
where $M_{ij} = a_{ii} a_{jj} - a_{ij}^{2}$ keeps rows and columns $i$ and $j$.
\begin{itemize}[leftmargin=1.6em,itemsep=0.3ex]
    \item psd $\Leftrightarrow$ all seven are $\ge 0$;
    \item nsd $\Leftrightarrow$ the three of order $1$ are $\le 0$, the three of
          order $2$ are $\ge 0$, and $\det A \le 0$;
    \item neither $\Rightarrow$ indefinite.
\end{itemize}
A good order of work: scan the diagonal (shortcuts), then the leading
minors $D_{1}, D_{2}, D_{3}$; only if one of them is $0$, check all seven.
\end{rulebox}

\begin{rulebox}{The $2 \times 2$ case}
For $A = \bigl(\begin{smallmatrix} a & b \\ b & c \end{smallmatrix}\bigr)$:
\begin{itemize}[leftmargin=1.6em,itemsep=0.3ex]
    \item pd $\Leftrightarrow$ $a > 0$ and $a c - b^{2} > 0$;
    \item nd $\Leftrightarrow$ $a < 0$ and $a c - b^{2} > 0$;
    \item indefinite $\Leftrightarrow$ $a c - b^{2} < 0$.
\end{itemize}
\end{rulebox}

\begin{rulebox}{Application: Hessian and second-order conditions}
For $f : \R^{n} \to \R$ of class $C^{2}$ at an interior critical point
$x^{*}$ (i.e.\ $\grad f(x^{*}) = 0$):
\begin{itemize}[leftmargin=1.6em,itemsep=0.3ex]
    \item Hessian $H_{f}(x^{*})$ pd $\Rightarrow$ strict local minimum.
    \item Hessian nd $\Rightarrow$ strict local maximum.
    \item Hessian indefinite $\Rightarrow$ saddle point.
    \item Otherwise (semidefinite, not definite) the test is inconclusive.
\end{itemize}
\end{rulebox}

\section{Worked Examples}

\begin{worked}{Write $Q(x,y)=2x^{2}-2xy+2y^{2}$ in matrix form and classify it by both criteria; then classify $x^{2}+4xy+y^{2}$ and $x^{2}+2xy+y^{2}$}
\wstep{Step 1. Build the matrix --- and halve the cross terms.}
A quadratic form $Q(\mathbf x)=\mathbf x^{\!\top}A\mathbf x$ with $A$
symmetric has $a_{ii}$ equal to the coefficient of $x_{i}^{2}$ and
$a_{ij}=a_{ji}$ equal to \textbf{half} the coefficient of $x_{i}x_{j}$. The
cross term $-2xy$ therefore contributes $-1$ in each off-diagonal slot:
\[
  Q(x,y)=2x^{2}-2xy+2y^{2}
  =\begin{pmatrix}x&y\end{pmatrix}
   \begin{pmatrix}2&-1\\-1&2\end{pmatrix}
   \begin{pmatrix}x\\y\end{pmatrix}.
\]
\wtrap{This halving is the chapter's central trap. Writing
$\begin{pmatrix}2&-2\\0&2\end{pmatrix}$ or
$\begin{pmatrix}2&-2\\-2&2\end{pmatrix}$ gives a different form and a
different classification. Check by expanding back: $\mathbf x^{\!\top}A
\mathbf x$ collects $a_{12}xy$ \emph{and} $a_{21}yx$, so the two halves
reassemble the full $-2xy$.}

\wstep{Step 2. First criterion: the eigenvalues.}
\[
  \det(A-\lambda I)=(2-\lambda)^{2}-1=\lambda^{2}-4\lambda+3=(\lambda-1)(\lambda-3),
\]
so $\lambda=1$ and $\lambda=3$. Both are strictly positive, hence
\[
  \boxed{A\ \text{is positive definite},\qquad Q(x,y)>0\ \text{for all}\ (x,y)\neq(0,0).}
\]

\wstep{Step 3. Second criterion: the leading principal minors (Sylvester).}
\[
  D_{1}=a_{11}=2>0,
  \qquad
  D_{2}=\det A=4-1=3>0 .
\]
All leading principal minors positive $\Rightarrow$ positive definite:
the same verdict, reached without solving a polynomial. For negative
definiteness the signs must \textbf{alternate} starting negative:
$D_{1}<0$, $D_{2}>0$, $D_{3}<0,\dots$
\wtrap{``All minors negative'' is not the test for negative definiteness ---
$D_{2}$ must be \emph{positive}. The reason is that $-A$ is positive
definite exactly when the minors of $A$ alternate, and $\det(-A)=(-1)^{n}
\det A$ flips every other sign.}

\wstep{Step 4. Confirm by completing the square, which shows \emph{why}.}
\[
  2x^{2}-2xy+2y^{2}
  =2\Bigl(x-\tfrac{y}{2}\Bigr)^{2}+\tfrac32 y^{2} .
\]
A sum of two positive multiples of squares, zero only when $y=0$ and then
$x=0$. Three independent routes, one answer.

\wstep{Step 5. Second form: $x^{2}+4xy+y^{2}$.}
Halving the cross term gives
$A=\begin{pmatrix}1&2\\2&1\end{pmatrix}$, and
\[
  D_{1}=1>0,\qquad D_{2}=1-4=-3<0 .
\]
A negative determinant in the $2\times2$ case means the eigenvalues have
opposite signs, so the form is \textbf{indefinite}: positive in some
directions, negative in others. Witness it directly --- $Q(1,0)=1>0$ and
$Q(1,-1)=1-4+1=-2<0$.

\wstep{Step 6. Third form: $x^{2}+2xy+y^{2}$, and the limit of Sylvester's test.}
Here $A=\begin{pmatrix}1&1\\1&1\end{pmatrix}$, so
\[
  D_{1}=1>0,\qquad D_{2}=1-1=0 .
\]
The leading-minor test is now \textbf{inconclusive}, and the answer is not
``positive definite''. In fact $Q(x,y)=(x+y)^{2}\ge0$ always, but it
\emph{vanishes} on the whole line $y=-x$: the form is positive
\textbf{semi}definite, not definite.
\wtrap{A zero leading minor ends Sylvester's test. Semidefiniteness requires
checking \emph{all} principal minors --- every diagonal sub-block, not only
the top-left ones --- and here they are $1$, $1$ and $0$, all $\ge0$, which
confirms positive semidefiniteness. Reporting ``positive definite'' from
$D_{1}>0,\ D_{2}=0$ is the error, and in \volref{optimization} it turns a saddle or a
flat valley into a claimed strict minimum.}

\wstep{Step 7. Why any of this matters.}
Replace $A$ by the Hessian of $f$ at a critical point:
\begin{center}\footnotesize
\begin{tabular}{ll}
  \hline
  $\Hess f$ at a critical point & conclusion\\
  \hline
  positive definite & strict local minimum\\
  negative definite & strict local maximum\\
  indefinite & saddle point\\
  semidefinite & the test says nothing\\
  \hline
\end{tabular}
\end{center}
The three forms above are, in order, a valley, a saddle, and a flat-bottomed
trough --- which is exactly the trio of pictures the second-order conditions
have to distinguish.
\end{worked}

\begin{worked}{Classify $A=\begin{pmatrix}-2&1&1\\1&-2&1\\1&1&-2\end{pmatrix}$, naming the rule used at each stage}
\wstep{Step 1. Scan the diagonal first.}
All three diagonal entries equal $-2<0$. No shortcut applies: there are no
opposite signs and no zero on the diagonal. The diagonal only tells us
what is still possible, namely nd or nsd (or indefinite); pd and psd are
already ruled out, since $Q(e_{1})=a_{11}=-2<0$.

\wstep{Step 2. Leading minors (Sylvester).}
\[
  D_{1}=-2<0,\qquad
  D_{2}=\begin{vmatrix}-2&1\\1&-2\end{vmatrix}=4-1=3>0,
\]
and, expanding along the first row,
\[
  D_{3}=\det A=-2\,(4-1)-1\,(-2-1)+1\,(1+2)=-6+3+3=0 .
\]
The pattern starts $-,+$ as negative definiteness requires, but
$D_{3}=0$, not $<0$. So $A$ is \textbf{not} negative definite, and since a
leading minor vanishes the test says nothing more.
\wtrap{Do not report ``indefinite'' because the pattern $-,+,0$ is not the
alternating one. The rule for other patterns applies only when all
leading minors are non-zero. A zero minor sends us to the full check.}

\wstep{Step 3. All seven principal minors.}
Order $1$: $a_{11}=a_{22}=a_{33}=-2\le0$. Order $2$: by symmetry
$M_{12}=M_{13}=M_{23}=(-2)(-2)-1^{2}=3\ge0$. Order $3$: $\det A=0\le0$.
Odd orders $\le0$, even orders $\ge0$: the matrix is
\[
  \boxed{\text{negative semidefinite, not negative definite.}}
\]

\wstep{Step 4. Confirm: a zero direction and a sum of squares.}
Every row sums to $0$, so $A\,(1,1,1)^{T}=0$ and $Q(1,1,1)=0$ with
$(1,1,1)\neq0$: definiteness fails, as the zero determinant said. Writing
out the form,
\[
  Q=-2x^{2}-2y^{2}-2z^{2}+2xy+2xz+2yz
   =-\bigl[(x-y)^{2}+(y-z)^{2}+(z-x)^{2}\bigr]\le0 ,
\]
which shows the semidefiniteness directly. (The eigenvalues are $0,-3,-3$.)
\end{worked}

\section{Practice Problems}

\begin{problems}
    \item Classify $A = \bigl(\begin{smallmatrix} 2 & 0 \\ 0 & 3 \end{smallmatrix}\bigr)$.
    \item Classify $A = \bigl(\begin{smallmatrix} -1 & 0 \\ 0 & -4 \end{smallmatrix}\bigr)$.
    \item Classify $A = \bigl(\begin{smallmatrix} 1 & 0 \\ 0 & -1 \end{smallmatrix}\bigr)$.
    \item Classify $A = \bigl(\begin{smallmatrix} 2 & 1 \\ 1 & 2 \end{smallmatrix}\bigr)$.
    \item Classify $A = \bigl(\begin{smallmatrix} 1 & 2 \\ 2 & 1 \end{smallmatrix}\bigr)$.
    \item Classify $A = \bigl(\begin{smallmatrix} 4 & 2 \\ 2 & 1 \end{smallmatrix}\bigr)$.
    \item For which $a$ is $A = \bigl(\begin{smallmatrix} a & 1 \\ 1 & a \end{smallmatrix}\bigr)$ positive definite?
    \item For which $b$ is $A = \bigl(\begin{smallmatrix} 4 & b \\ b & 9 \end{smallmatrix}\bigr)$ positive definite?
    \item Write $Q(x, y) = 3 x^{2} + 4 x y + y^{2}$ as $x^{T} A x$ for a symmetric $A$, and classify $A$.
    \item Write $Q(x, y) = x^{2} - 6 x y + 10 y^{2}$ as $x^{T} A x$ and classify.
    \item Classify $A = \begin{pmatrix} 2 & 1 & 0 \\ 1 & 2 & 1 \\ 0 & 1 & 2 \end{pmatrix}$ via leading principal minors.
    \item Classify $A = \begin{pmatrix} -2 & 1 & 0 \\ 1 & -3 & 1 \\ 0 & 1 & -2 \end{pmatrix}$ via leading principal minors.
    \item Classify $A = \begin{pmatrix} 1 & 1 & 0 \\ 1 & 2 & 1 \\ 0 & 1 & 3 \end{pmatrix}$.
    \item Classify $A = \begin{pmatrix} 1 & 2 & 0 \\ 2 & 1 & 0 \\ 0 & 0 & 5 \end{pmatrix}$.
    \item For $f(x, y) = x^{2} + x y + y^{2}$, find the Hessian and classify the critical point at the origin.
    \item For $f(x, y) = x^{2} - x y + y^{2} - 3 x$, find the critical point and classify it.
    \item For $f(x, y) = x^{3} - 3 x + y^{2}$, find all critical points and classify them.
    \item For $f(x, y) = x^{2} + y^{2} - x y - x - y$, find the critical point and classify.
    \item Show that $A^{T} A$ is psd for any real matrix $A$.
    \item Show that if $A$ is pd then $A$ is invertible.
    \item Find the eigenvalues of $A = \bigl(\begin{smallmatrix} 3 & -1 \\ -1 & 3 \end{smallmatrix}\bigr)$ and classify it.
    \item A covariance matrix is symmetric and psd. If $\Sigma = \bigl(\begin{smallmatrix} 4 & \rho \\ \rho & 1 \end{smallmatrix}\bigr)$ is a valid covariance matrix, what range of $\rho$ is allowed?
    \item For the same $\Sigma$ with $\rho = 1$, is $\Sigma$ pd or just psd? Justify.
    \item Compute $Q(1, 2) = x^{T} A x$ for $A = \bigl(\begin{smallmatrix} 1 & 2 \\ 2 & 3 \end{smallmatrix}\bigr)$.
    \item For which $\alpha$ is the Hessian of $f(x, y) = \alpha x^{2} - 2 x y + y^{2}$ pd?
    \item Classify $A = \begin{pmatrix} 3 & 1 & -1 \\ 1 & 2 & 0 \\ -1 & 0 & 4 \end{pmatrix}$ and name the rule you use.
    \item Classify $A = \begin{pmatrix} -3 & 1 & 1 \\ 1 & -2 & 0 \\ 1 & 0 & -4 \end{pmatrix}$ and name the rule you use.
    \item Classify $A = \begin{pmatrix} -1 & 0 & 2 \\ 0 & -1 & 2 \\ 2 & 2 & -1 \end{pmatrix}$ and name the rule you use. Then exhibit vectors on which $Q$ takes both signs.
    \item Classify $A = \begin{pmatrix} 1 & 3 & 0 \\ 3 & 2 & 1 \\ 0 & 1 & 5 \end{pmatrix}$ with as little computation as possible.
    \item Classify $A = \begin{pmatrix} 4 & 1 & 2 \\ 1 & -1 & 0 \\ 2 & 0 & 3 \end{pmatrix}$ without computing a single determinant.
    \item Classify $A = \begin{pmatrix} 1 & -1 & 2 \\ -1 & 5 & -2 \\ 2 & -2 & 4 \end{pmatrix}$. If it is semidefinite, find a non-zero $x$ with $x^{T} A x = 0$.
    \item Classify $A = \begin{pmatrix} -1 & 2 & -1 \\ 2 & -5 & 1 \\ -1 & 1 & -2 \end{pmatrix}$. If it is semidefinite, find a non-zero $x$ with $x^{T} A x = 0$.
\end{problems}

\section{Solutions}

\begin{solutions}
    \item Eigenvalues $2, 3 > 0$: positive definite.
    \item Eigenvalues $-1, -4 < 0$: negative definite.
    \item Eigenvalues $1, -1$: indefinite.
    \item Eigenvalues $1, 3 > 0$: pd. (Or: $D_{1} = 2 > 0$, $D_{2} = 3 > 0$.)
    \item $D_{1} = 1 > 0$, $D_{2} = 1 - 4 = -3 < 0$: indefinite. (Eigenvalues $-1, 3$.)
    \item $D_{2} = 4 - 4 = 0$: not pd. Eigenvalues from $\lambda^{2} - 5 \lambda + 0 = 0$ give $\lambda = 0, 5$. Both $\ge 0$ with one zero: psd but not pd.
    \item $D_{1} = a > 0$ and $D_{2} = a^{2} - 1 > 0$, i.e.\ $|a| > 1$. Combined: $a > 1$.
    \item $D_{1} = 4 > 0$, $D_{2} = 36 - b^{2} > 0 \Leftrightarrow |b| < 6$.
    \item $A = \bigl(\begin{smallmatrix} 3 & 2 \\ 2 & 1 \end{smallmatrix}\bigr)$. $D_{1} = 3 > 0$, $D_{2} = 3 - 4 = -1 < 0$: indefinite.
    \item $A = \bigl(\begin{smallmatrix} 1 & -3 \\ -3 & 10 \end{smallmatrix}\bigr)$. $D_{1} = 1 > 0$, $D_{2} = 10 - 9 = 1 > 0$: pd.
    \item $D_{1} = 2$, $D_{2} = 4 - 1 = 3$, $D_{3} = $ expand along row 1: $2(4 - 1) - 1(2 - 0) + 0 = 6 - 2 = 4$. All $> 0$: pd.
    \item $D_{1} = -2$, $D_{2} = 6 - 1 = 5$, $D_{3} = $ expand: $-2(6 - 1) - 1(-2 - 0) + 0 = -10 + 2 = -8$. Signs $-, +, -$: nd.
    \item $D_{1} = 1$, $D_{2} = 2 - 1 = 1$, $D_{3}$ by expansion: $1(6 - 1) - 1(3 - 0) + 0 = 5 - 3 = 2$. All $> 0$: pd.
    \item Block diagonal. Upper $2 \times 2$ has $D_{1} = 1 > 0$, $D_{2} = 1 - 4 = -3 < 0$: indefinite. Thus $A$ has at least one negative and one positive eigenvalue (from the block), so $A$ is indefinite overall.
    \item $f_{x} = 2 x + y$, $f_{y} = x + 2 y$. Hessian $H = \bigl(\begin{smallmatrix} 2 & 1 \\ 1 & 2 \end{smallmatrix}\bigr)$. $D_{1} = 2 > 0$, $D_{2} = 3 > 0$: pd. Origin is a strict local (and global) min.
    \item $f_{x} = 2 x - y - 3 = 0$, $f_{y} = -x + 2 y = 0 \Rightarrow x = 2 y$. Substituting: $4 y - y - 3 = 0 \Rightarrow y = 1, x = 2$. Hessian $H = \bigl(\begin{smallmatrix} 2 & -1 \\ -1 & 2 \end{smallmatrix}\bigr)$, $D_{1} = 2 > 0$, $D_{2} = 3 > 0$: pd. Strict local min at $(2, 1)$.
    \item $f_{x} = 3 x^{2} - 3 = 0 \Rightarrow x = \pm 1$. $f_{y} = 2 y = 0 \Rightarrow y = 0$. Critical points $(1, 0)$ and $(-1, 0)$. $H = \bigl(\begin{smallmatrix} 6 x & 0 \\ 0 & 2 \end{smallmatrix}\bigr)$. At $(1, 0)$: $H = \bigl(\begin{smallmatrix} 6 & 0 \\ 0 & 2 \end{smallmatrix}\bigr)$, pd $\Rightarrow$ local min. At $(-1, 0)$: $H = \bigl(\begin{smallmatrix} -6 & 0 \\ 0 & 2 \end{smallmatrix}\bigr)$, indefinite $\Rightarrow$ saddle.
    \item $f_{x} = 2 x - y - 1 = 0$, $f_{y} = 2 y - x - 1 = 0$. Solve: from the first, $y = 2 x - 1$; substitute: $2(2 x - 1) - x - 1 = 3 x - 3 = 0$, so $x = 1, y = 1$. $H = \bigl(\begin{smallmatrix} 2 & -1 \\ -1 & 2 \end{smallmatrix}\bigr)$, pd. Local min at $(1, 1)$.
    \item $x^{T} (A^{T} A) x = (A x)^{T} (A x) = \|A x\|^{2} \ge 0$ for all $x$, so $A^{T} A$ is psd.
    \item If $A$ is pd then no eigenvalue is $0$ (all $> 0$), so $\det A = \prod \lambda_{i} > 0$, hence $A$ is invertible.
    \item $p(\lambda) = (3 - \lambda)^{2} - 1 = \lambda^{2} - 6 \lambda + 8 = (\lambda - 2)(\lambda - 4)$. Eigenvalues $2, 4 > 0$: pd.
    \item $\Sigma$ psd $\Leftrightarrow$ $D_{1} = 4 \ge 0$ (yes) and $D_{2} = 4 - \rho^{2} \ge 0 \Leftrightarrow |\rho| \le 2$. (Note: a valid \emph{correlation} matrix would force $|\rho| \le 1$ since correlations lie in $[-1, 1]$; the question only asks psd-ness for the covariance.)
    \item With $\rho = 1$: $D_{2} = 4 - 1 = 3 > 0$, so $\Sigma$ is pd (eigenvalues are roots of $\lambda^{2} - 5\lambda + 3 = 0$, both positive).
    \item $x^{T} A x = (1, 2)\bigl(\begin{smallmatrix} 1 & 2 \\ 2 & 3 \end{smallmatrix}\bigr)\bigl(\begin{smallmatrix} 1 \\ 2 \end{smallmatrix}\bigr) = (1, 2)\bigl(\begin{smallmatrix} 5 \\ 8 \end{smallmatrix}\bigr) = 5 + 16 = 21$.
    \item Hessian $H = \bigl(\begin{smallmatrix} 2 \alpha & -2 \\ -2 & 2 \end{smallmatrix}\bigr)$. pd $\Leftrightarrow$ $2 \alpha > 0$ and $4 \alpha - 4 > 0$, i.e.\ $\alpha > 1$.
    \item Diagonal $3, 2, 4$: all positive, so pd is still possible (no shortcut decides). Leading minors (Sylvester): $D_{1} = 3$, $D_{2} = 3 \cdot 2 - 1^{2} = 5$, and expanding along row 1,
          $D_{3} = 3\,(2 \cdot 4 - 0) - 1\,(1 \cdot 4 - 0 \cdot (-1)) + (-1)\,(1 \cdot 0 - 2 \cdot (-1)) = 24 - 4 - 2 = 18$.
          All three $> 0$: positive definite by Sylvester's criterion.
    \item Diagonal $-3, -2, -4$: all negative, so nd is still possible. Leading minors: $D_{1} = -3$, $D_{2} = 6 - 1 = 5$, and
          $D_{3} = -3\,(8 - 0) - 1\,(-4 - 0) + 1\,(0 - (-2)) = -24 + 4 + 2 = -18$.
          Signs $-, +, -$ alternate starting negative: negative definite by Sylvester's criterion.
    \item Diagonal $-1, -1, -1$: no shortcut. Leading minors: $D_{1} = -1$, $D_{2} = 1 - 0 = 1$, and
          $D_{3} = -1\,(1 - 4) - 0 + 2\,(0 \cdot 2 - (-1) \cdot 2) = 3 + 4 = 7$.
          All three are non-zero and the pattern $-, +, +$ is neither all positive nor alternating, so $A$ is indefinite (rule for the other sign patterns).
          Witnesses: $Q(1, 0, 0) = a_{11} = -1 < 0$, while $Q(1, 1, 1)$ is the sum of all nine entries, $-1 + 0 + 2 + 0 - 1 + 2 + 2 + 2 - 1 = 5 > 0$.
    \item Diagonal $1, 2, 5$: no shortcut. Leading minors: $D_{1} = 1 > 0$, $D_{2} = 1 \cdot 2 - 3^{2} = -7 < 0$. A negative minor of even order means indefinite; $D_{3}$ is not needed (it equals $-36$).
          Witnesses: $Q(1, 0, 0) = 1 > 0$ and $Q(3, -2, 0) = 9 - 36 + 8 = -19 < 0$.
    \item Diagonal shortcut: $a_{11} = 4 > 0$ and $a_{22} = -1 < 0$ have opposite signs, so $A$ is indefinite. Indeed $Q(e_{1}) = 4 > 0$ and $Q(e_{2}) = -1 < 0$. No determinant is needed.
    \item Diagonal $1, 5, 4$: no shortcut. Leading minors: $D_{1} = 1$, $D_{2} = 5 - 1 = 4$,
          $D_{3} = 1\,(20 - 4) - (-1)\,(-4 + 4) + 2\,(2 - 10) = 16 - 0 - 16 = 0$.
          A zero leading minor: not pd, and Sylvester stops. All seven principal minors: order 1: $1, 5, 4$; order 2: $M_{12} = 5 - 1 = 4$, $M_{13} = 4 - 4 = 0$, $M_{23} = 20 - 4 = 16$; order 3: $\det A = 0$. All seven $\ge 0$: positive semidefinite, not pd.
          Zero direction: solve $A x = 0$. The third row is twice the first, so only rows 1 and 2 count: $x_{1} - x_{2} + 2 x_{3} = 0$ and $-x_{1} + 5 x_{2} - 2 x_{3} = 0$. Adding them gives $4 x_{2} = 0$, so $x_{2} = 0$ and $x_{1} = -2 x_{3}$. With $x_{3} = -1$: $x = (2, 0, -1)$, and $Q(2, 0, -1) = 0$.
          (Check: $Q = (x - y + 2z)^{2} + 4 y^{2} \ge 0$.)
    \item Diagonal $-1, -5, -2$: no shortcut. Leading minors: $D_{1} = -1$, $D_{2} = 5 - 4 = 1$,
          $D_{3} = -1\,(10 - 1) - 2\,(-4 + 1) + (-1)\,(2 - 5) = -9 + 6 + 3 = 0$.
          Not nd (needs $D_{3} < 0$), and Sylvester stops. All seven principal minors: order 1: $-1, -5, -2$, all $\le 0$; order 2: $M_{12} = 5 - 4 = 1$, $M_{13} = 2 - 1 = 1$, $M_{23} = 10 - 1 = 9$, all $\ge 0$; order 3: $\det A = 0 \le 0$. Negative semidefinite, not nd.
          Zero direction: from rows 1 and 3, $-x_{1} + 2 x_{2} - x_{3} = 0$ and $-x_{1} + x_{2} - 2 x_{3} = 0$; subtracting gives $x_{2} = -x_{3}$, then $x_{1} = -3 x_{3}$. With $x_{3} = -1$: $x = (3, 1, -1)$, and row 2 checks: $6 - 5 - 1 = 0$. So $Q(3, 1, -1) = 0$.
          (Check: $Q = -(x - 2y + z)^{2} - (y + z)^{2} \le 0$.)
\end{solutions}

\section{Additional Problems}

\begin{problems}
    \item Classify $A = \bigl(\begin{smallmatrix} 5 & 0 \\ 0 & 1 \end{smallmatrix}\bigr)$.
    \item Classify $A = \bigl(\begin{smallmatrix} -3 & 0 \\ 0 & -2 \end{smallmatrix}\bigr)$.
    \item Classify $A = \bigl(\begin{smallmatrix} 3 & 0 \\ 0 & -1 \end{smallmatrix}\bigr)$.
    \item Classify $A = \bigl(\begin{smallmatrix} 3 & 1 \\ 1 & 3 \end{smallmatrix}\bigr)$.
    \item Classify $A = \bigl(\begin{smallmatrix} 2 & 3 \\ 3 & 2 \end{smallmatrix}\bigr)$.
    \item Classify $A = \bigl(\begin{smallmatrix} 9 & 6 \\ 6 & 4 \end{smallmatrix}\bigr)$.
    \item For which $b$ is $A = \bigl(\begin{smallmatrix} 3 & b \\ b & 3 \end{smallmatrix}\bigr)$ positive definite?
    \item For which $c$ is $A = \bigl(\begin{smallmatrix} -2 & 1 \\ 1 & c \end{smallmatrix}\bigr)$ negative definite?
    \item Write $Q(x, y) = 2 x^{2} + 6 x y + 5 y^{2}$ as $x^{T} A x$ and classify $A$.
    \item Write $Q(x, y) = -x^{2} + 4 x y - 5 y^{2}$ as $x^{T} A x$ and classify.
    \item Classify $A = \begin{pmatrix} 3 & 1 & 0 \\ 1 & 3 & 1 \\ 0 & 1 & 3 \end{pmatrix}$ via leading principal minors.
    \item Classify $A = \begin{pmatrix} -3 & 2 & 0 \\ 2 & -4 & 1 \\ 0 & 1 & -2 \end{pmatrix}$ via leading principal minors.
    \item Classify $A = \begin{pmatrix} 2 & 1 & 1 \\ 1 & 2 & 1 \\ 1 & 1 & 2 \end{pmatrix}$.
    \item Classify $A = \begin{pmatrix} 2 & 0 & 0 \\ 0 & -1 & 0 \\ 0 & 0 & 3 \end{pmatrix}$.
    \item For $f(x, y) = 2 x^{2} - 4 x y + 3 y^{2}$, find the Hessian and classify the origin.
    \item For $f(x, y) = x^{2} + 4 x y + 2 y^{2} + 6 x$, find the critical point and classify it.
    \item For $f(x, y) = x^{3} - 12 x + y^{2} - 2 y$, find all critical points and classify them.
    \item For $f(x, y) = x^{2} + 4 x y + y^{2} - 2 x - 4 y$, find the critical point and classify.
    \item Show that if $A$ is symmetric and all eigenvalues are positive, then $A$ is invertible and $A^{-1}$ has all positive eigenvalues.
    \item Show: $A$ is positive definite iff $-A$ is negative definite.
    \item Find the eigenvalues of $A = \bigl(\begin{smallmatrix} 4 & 1 \\ 1 & 4 \end{smallmatrix}\bigr)$ and classify it.
    \item For $\Sigma = \bigl(\begin{smallmatrix} 9 & \rho \\ \rho & 16 \end{smallmatrix}\bigr)$ a valid covariance matrix, what range of $\rho$ is allowed?
    \item Compute $Q(2, -1) = x^{T} A x$ for $A = \bigl(\begin{smallmatrix} 3 & 1 \\ 1 & 2 \end{smallmatrix}\bigr)$.
    \item For which $\alpha$ is the Hessian of $f(x, y) = x^{2} + 2 \alpha x y + 4 y^{2}$ pd?
    \item Determine whether $Q(x, y, z) = x^{2} + y^{2} + z^{2} + 2 x y$ is pd, psd, or indefinite.
    \item Classify $A = \begin{pmatrix} 2 & 1 & 3 \\ 1 & 0 & 0 \\ 3 & 0 & 4 \end{pmatrix}$ by a diagonal shortcut, and give a vector $x$ with $x^{T} A x < 0$.
    \item Classify $A = \begin{pmatrix} 1 & 2 & -1 \\ 2 & 4 & -2 \\ -1 & -2 & 1 \end{pmatrix}$ and write $Q(x, y, z)$ as a single square.
    \item Classify $A = \begin{pmatrix} -2 & 0 & 1 \\ 0 & 0 & 0 \\ 1 & 0 & -3 \end{pmatrix}$. Does the zero on the diagonal make it indefinite?
    \item Classify $A = \begin{pmatrix} 1 & 1 & 0 \\ 1 & 1 & 1 \\ 0 & 1 & 1 \end{pmatrix}$. Its second leading minor is $0$: what does the full check give?
    \item True or false (justify): a symmetric $3 \times 3$ matrix whose diagonal entries are all positive is positive definite.
    \item For which $t$ is $A = \begin{pmatrix} 1 & 0 & t \\ 0 & 2 & 0 \\ t & 0 & 4 \end{pmatrix}$ positive definite, positive semidefinite but not pd, or indefinite?
\end{problems}

\section{Additional Solutions}

\begin{solutions}
    \item Eigenvalues $5, 1 > 0$: pd.
    \item Eigenvalues $-3, -2 < 0$: nd.
    \item Eigenvalues $3, -1$: indefinite.
    \item Eigenvalues from $p(\lambda) = (3 - \lambda)^{2} - 1 = (\lambda - 2)(\lambda - 4)$: both positive, pd. Or: $D_{1} = 3 > 0$, $D_{2} = 9 - 1 = 8 > 0$. pd.
    \item $D_{1} = 2 > 0$, $D_{2} = 4 - 9 = -5 < 0$: indefinite. (Eigenvalues $-1, 5$.)
    \item $D_{2} = 36 - 36 = 0$: not pd. Eigenvalues from $\lambda^{2} - 13 \lambda + 0 = 0$: $\lambda = 0, 13$. psd, not pd.
    \item $D_{1} = 3 > 0$, $D_{2} = 9 - b^{2} > 0 \Leftrightarrow |b| < 3$.
    \item nd requires $D_{1} = -2 < 0$ (OK) and $D_{2} = -2 c - 1 > 0 \Leftrightarrow c < -1/2$.
    \item $A = \bigl(\begin{smallmatrix} 2 & 3 \\ 3 & 5 \end{smallmatrix}\bigr)$. $D_{1} = 2 > 0$, $D_{2} = 10 - 9 = 1 > 0$: pd.
    \item $A = \bigl(\begin{smallmatrix} -1 & 2 \\ 2 & -5 \end{smallmatrix}\bigr)$. $D_{1} = -1 < 0$, $D_{2} = 5 - 4 = 1 > 0$: signs $-, +$: nd.
    \item $D_{1} = 3, D_{2} = 9 - 1 = 8, D_{3} = $ expand along row 1: $3(9 - 1) - 1(3 - 0) + 0 = 24 - 3 = 21$. All $> 0$: pd.
    \item $D_{1} = -3, D_{2} = 12 - 4 = 8, D_{3} = $ expand: $-3(8 - 1) - 2(-4 - 0) + 0 = -21 + 8 = -13$. Signs $-, +, -$: nd.
    \item $D_{1} = 2, D_{2} = 4 - 1 = 3, D_{3} = $ expand: $2(4 - 1) - 1(2 - 1) + 1(1 - 2) = 6 - 1 - 1 = 4$. All $> 0$: pd. (Eigenvalues are $1, 1, 4$.)
    \item Diagonal: eigenvalues $2, -1, 3$ (mixed signs): indefinite.
    \item $H = \bigl(\begin{smallmatrix} 4 & -4 \\ -4 & 6 \end{smallmatrix}\bigr)$. $D_{1} = 4 > 0$, $D_{2} = 24 - 16 = 8 > 0$: pd. Origin is a strict local min.
    \item $f_{x} = 2 x + 4 y + 6 = 0$, $f_{y} = 4 x + 4 y = 0 \Rightarrow x = -y$. Substitute: $-2 y + 4 y + 6 = 0 \Rightarrow y = -3, x = 3$. $H = \bigl(\begin{smallmatrix} 2 & 4 \\ 4 & 4 \end{smallmatrix}\bigr)$, $D_{2} = 8 - 16 = -8 < 0$: indefinite. Saddle at $(3, -3)$.
    \item $f_{x} = 3 x^{2} - 12 = 0 \Rightarrow x = \pm 2$. $f_{y} = 2 y - 2 = 0 \Rightarrow y = 1$. Critical points $(2, 1)$ and $(-2, 1)$. $H = \bigl(\begin{smallmatrix} 6 x & 0 \\ 0 & 2 \end{smallmatrix}\bigr)$. At $(2, 1)$: $\bigl(\begin{smallmatrix} 12 & 0 \\ 0 & 2 \end{smallmatrix}\bigr)$, pd: local min. At $(-2, 1)$: $\bigl(\begin{smallmatrix} -12 & 0 \\ 0 & 2 \end{smallmatrix}\bigr)$, indefinite: saddle.
    \item $f_{x} = 2 x + 4 y - 2 = 0$, $f_{y} = 4 x + 2 y - 4 = 0$. Solve: $x + 2 y = 1$, $2 x + y = 2$. From first $x = 1 - 2 y$; sub: $2 - 4 y + y = 2 \Rightarrow y = 0, x = 1$. $H = \bigl(\begin{smallmatrix} 2 & 4 \\ 4 & 2 \end{smallmatrix}\bigr)$, $D_{2} = 4 - 16 = -12 < 0$: indefinite. Saddle at $(1, 0)$.
    \item All eigenvalues $> 0$ implies $\det A = \prod \lambda_{i} > 0$, so $A$ is invertible. The eigenvalues of $A^{-1}$ are $1/\lambda_{i}$, each $> 0$.
    \item $x^{T}(-A)x = -(x^{T} A x)$. If $x^{T} A x > 0$ for all $x \neq 0$, then $x^{T}(-A)x < 0$ for all $x \neq 0$. Converse identical.
    \item $p(\lambda) = (4 - \lambda)^{2} - 1 = (\lambda - 3)(\lambda - 5)$. Eigenvalues $3, 5 > 0$: pd.
    \item Need $D_{1} = 9 \ge 0$ (OK), $D_{2} = 144 - \rho^{2} \ge 0 \Leftrightarrow |\rho| \le 12$.
    \item $A \bigl(\begin{smallmatrix} 2 \\ -1 \end{smallmatrix}\bigr) = \bigl(\begin{smallmatrix} 6 - 1 \\ 2 - 2 \end{smallmatrix}\bigr) = \bigl(\begin{smallmatrix} 5 \\ 0 \end{smallmatrix}\bigr)$. $(2, -1) \cdot (5, 0) = 10$. So $Q = 10$.
    \item $H = \bigl(\begin{smallmatrix} 2 & 2 \alpha \\ 2 \alpha & 8 \end{smallmatrix}\bigr)$. pd iff $2 > 0$ (OK) and $16 - 4 \alpha^{2} > 0 \Leftrightarrow |\alpha| < 2$.
    \item $A = \begin{pmatrix} 1 & 1 & 0 \\ 1 & 1 & 0 \\ 0 & 0 & 1 \end{pmatrix}$. Eigenvalues: from $2 \times 2$ block $\bigl(\begin{smallmatrix} 1 & 1 \\ 1 & 1 \end{smallmatrix}\bigr)$ we get $0, 2$; from $1 \times 1$ block, $1$. So eigenvalues $\{0, 1, 2\}$. All $\ge 0$, one zero: psd but not pd. Verify: $Q = (x + y)^{2} + z^{2} \ge 0$, with $Q = 0$ at $(1, -1, 0) \neq 0$.
    \item Diagonal shortcut: $a_{22} = 0$ while $a_{21} = 1 \neq 0$ in the same row, so $A$ is indefinite. The witness is the $2 \times 2$ principal minor on rows and columns $1, 2$: $M_{12} = 2 \cdot 0 - 1^{2} = -1 < 0$.
          On the plane $z = 0$ the form is $Q(x, y, 0) = 2 x^{2} + 2 x y$. At $(1, 0, 0)$ it is $2 > 0$; at $(1, -2, 0)$ it is $2 - 4 = -2 < 0$.
    \item Diagonal $1, 4, 1$: no shortcut. Leading minors: $D_{1} = 1$, $D_{2} = 4 - 4 = 0$, so Sylvester stops (and pd is excluded). All seven principal minors: order 1: $1, 4, 1$; order 2: $M_{12} = 4 - 4 = 0$, $M_{13} = 1 - 1 = 0$, $M_{23} = 4 - 4 = 0$; order 3: $\det A = 0$ (all rows are multiples of $(1, 2, -1)$). All seven $\ge 0$: positive semidefinite, not pd.
          Single square: $Q = x^{2} + 4 y^{2} + z^{2} + 4 x y - 2 x z - 4 y z = (x + 2 y - z)^{2}$, which vanishes on the whole plane $x + 2y - z = 0$, for example at $(2, -1, 0)$.
    \item The zero diagonal entry $a_{22} = 0$ sits in a row that is entirely zero, so the shortcut does \emph{not} apply. Leading minors: $D_{1} = -2$, $D_{2} = 0$, $D_{3} = 0$ (zero row); Sylvester stops. All seven principal minors: order 1: $-2, 0, -3$, all $\le 0$; order 2: $M_{12} = 0$, $M_{13} = 6 - 1 = 5$, $M_{23} = 0$, all $\ge 0$; order 3: $\det A = 0 \le 0$. Negative semidefinite, not nd.
          Indeed $Q = -2 x^{2} + 2 x z - 3 z^{2}$ does not involve $y$; it is negative definite in $(x, z)$ (since $-2 < 0$, $M_{13} = 5 > 0$) and zero along $(0, 1, 0)$.
    \item Diagonal $1, 1, 1$: no shortcut. Leading minors: $D_{1} = 1$, $D_{2} = 1 - 1 = 0$, $D_{3} = 1\,(1 - 1) - 1\,(1 - 0) + 0 = -1$. With a zero leading minor go to the full check: order 1: $1, 1, 1$; order 2: $M_{12} = 0$, $M_{13} = 1$, $M_{23} = 0$; order 3: $\det A = -1$.
          Not psd ($\det A < 0$) and not nsd (diagonal positive): indefinite.
          Witnesses: $Q(1, 0, 0) = 1 > 0$ and $Q(1, -1, 1) = 1 + 1 + 1 - 2 - 2 + 0 = -1 < 0$.
    \item False. Positive diagonal entries are necessary for pd but not sufficient. Counterexample: $A = \begin{pmatrix} 2 & 0 & 3 \\ 0 & 1 & 0 \\ 3 & 0 & 2 \end{pmatrix}$ has diagonal $2, 1, 2$, yet the principal minor $M_{13} = 4 - 9 = -5 < 0$, so it is indefinite: $Q(1, 0, -1) = 2 + 2 - 6 = -2 < 0$.
    \item Principal minors: order 1: $1, 2, 4$, all positive; order 2: $M_{12} = 2$, $M_{13} = 4 - t^{2}$, $M_{23} = 8$; order 3, expanding along the middle row: $\det A = 2\,(4 - t^{2})$. The leading minors are $D_{1} = 1$, $D_{2} = 2$, $D_{3} = 2\,(4 - t^{2})$.
          pd (Sylvester): $D_{3} > 0 \Leftrightarrow |t| < 2$.
          psd but not pd: all seven $\ge 0$ needs $4 - t^{2} \ge 0$, and not pd forces $4 - t^{2} = 0$, so $t = \pm 2$. (At $t = 2$: $Q = (x + 2z)^{2} + 2 y^{2}$.)
          Indefinite: $|t| > 2$, where $M_{13} < 0$ (a negative minor of even order).
\end{solutions}
